\documentclass[11pt]{article}

\usepackage{microtype}
\usepackage{graphicx}
\usepackage{subcaption}
\usepackage{booktabs}
\usepackage{microsoft-tech-report}
\usepackage{hyperref}
\usepackage{amsmath}
\usepackage{amssymb}
\usepackage{mathtools}
\usepackage{amsthm}
\usepackage{bm}
\usepackage{enumitem}
\usepackage[capitalize,noabbrev]{cleveref}
\usepackage{xspace}
\usepackage{titletoc}
\usepackage{placeins}  % provides \FloatBarrier
\usepackage{tcolorbox} % colored boxes for theorems
\usepackage[ruled,linesnumbered,noend,noline]{algorithm2e} % pseudocode
\DontPrintSemicolon
\usepackage{listings}  % for code/JSON in prompt boxes
\tcbuselibrary{theorems,skins,breakable,listings}

% --- Prompt box for LLM prompts ---
\definecolor{promptbg}{HTML}{F7F7F8}
\definecolor{promptborder}{HTML}{5B5B5B}
\newtcolorbox{promptbox}[1][]{
  colback=promptbg, colframe=promptborder, fonttitle=\bfseries\sffamily,
  breakable, enhanced, arc=3pt, boxrule=0.8pt,
  left=10pt, right=10pt, top=8pt, bottom=8pt,
  title={#1}
}

% --- JSON listing style ---
\lstdefinestyle{json}{
  basicstyle=\small\ttfamily,
  breaklines=true,
  breakatwhitespace=false,
  showstringspaces=false,
  columns=fullflexible,
  literate={"}{\textquotedbl}1,
}

% --- Appendix TOC toggle ---
% Set to true to show a local table of contents at the beginning of the Appendix
\newif\ifappendixtoc
\appendixtoctrue   % <-- change to \appendixtocfalse to disable

\newcommand{\theHalgorithm}{\arabic{algorithm}}
\definecolor{color_blue}{HTML}{E7EFFA}
\definecolor{color_green}{HTML}{E6F8E0}
\definecolor{color_gray}{HTML}{ECECEC}
\definecolor{pearDark}{HTML}{2980B9}
\definecolor{theoremblue}{HTML}{EBF5FB}
\definecolor{theoremborder}{HTML}{2980B9}
\definecolor{propgreen}{HTML}{EAFAF1}
\definecolor{propborder}{HTML}{27AE60}
\definecolor{defyellow}{HTML}{FEF9E7}
\definecolor{defborder}{HTML}{F39C12}
\definecolor{remarkgray}{HTML}{F2F3F4}
\definecolor{remarkborder}{HTML}{7F8C8D}

\hypersetup{
  colorlinks=true,
  linkcolor=msftblue,
  citecolor=msftblue,
  urlcolor=msftblue,
  pdftitle={[Your Paper Title]}          % <-- Replace with your title
}

% --- Colored theorem-like environments ---
% Note: tcolorbox labels use the format {prefix}:{label}, e.g., \cref{thm:example}
\newtcbtheorem[auto counter]{theorem}{Theorem}{
  colback=theoremblue, colframe=theoremborder, fonttitle=\bfseries,
  breakable, enhanced, sharp corners, boxrule=0.6pt, left=6pt, right=6pt, top=6pt, bottom=6pt,
  crefname={Theorem}{Theorems}
}{thm}

\newtcbtheorem[use counter from=theorem]{proposition}{Proposition}{
  colback=propgreen, colframe=propborder, fonttitle=\bfseries,
  breakable, enhanced, sharp corners, boxrule=0.6pt, left=6pt, right=6pt, top=6pt, bottom=6pt,
  crefname={Proposition}{Propositions}
}{prop}

\newtcbtheorem[use counter from=theorem]{lemma}{Lemma}{
  colback=theoremblue, colframe=theoremborder, fonttitle=\bfseries,
  breakable, enhanced, sharp corners, boxrule=0.6pt, left=6pt, right=6pt, top=6pt, bottom=6pt,
  crefname={Lemma}{Lemmas}
}{lem}

\newtcbtheorem[use counter from=theorem]{corollary}{Corollary}{
  colback=propgreen, colframe=propborder, fonttitle=\bfseries,
  breakable, enhanced, sharp corners, boxrule=0.6pt, left=6pt, right=6pt, top=6pt, bottom=6pt,
  crefname={Corollary}{Corollaries}
}{cor}

\newtcbtheorem[use counter from=theorem]{definition}{Definition}{
  colback=defyellow, colframe=defborder, fonttitle=\bfseries,
  breakable, enhanced, sharp corners, boxrule=0.6pt, left=6pt, right=6pt, top=6pt, bottom=6pt,
  crefname={Definition}{Definitions}
}{def}

\newtcbtheorem[use counter from=theorem]{assumption}{Assumption}{
  colback=defyellow, colframe=defborder, fonttitle=\bfseries,
  breakable, enhanced, sharp corners, boxrule=0.6pt, left=6pt, right=6pt, top=6pt, bottom=6pt,
  crefname={Assumption}{Assumptions}
}{asm}

\newtcbtheorem[use counter from=theorem]{remark}{Remark}{
  colback=remarkgray, colframe=remarkborder, fonttitle=\bfseries,
  breakable, enhanced, sharp corners, boxrule=0.6pt, left=6pt, right=6pt, top=6pt, bottom=6pt,
  crefname={Remark}{Remarks}
}{rem}

% --- Register tcolorbox counters with cleveref ---
\crefname{tcb@cnt@theorem}{Theorem}{Theorems}
\crefname{tcb@cnt@proposition}{Proposition}{Propositions}
\crefname{tcb@cnt@lemma}{Lemma}{Lemmas}
\crefname{tcb@cnt@corollary}{Corollary}{Corollaries}
\crefname{tcb@cnt@definition}{Definition}{Definitions}
\crefname{tcb@cnt@assumption}{Assumption}{Assumptions}
\crefname{tcb@cnt@remark}{Remark}{Remarks}
\crefname{algocf}{Algorithm}{Algorithms}

% --- Custom commands (add your own below) ---
\newcommand{\method}{YourMethod\xspace}   % <-- Replace with your method name
\techreportshorttitle{[Short Title]}      % <-- Replace with your short title

\begin{document}
\thispagestyle{empty}

% ============================================================
%  Header: Microsoft logo + date
% ============================================================
\noindent
\begin{minipage}[c]{0.5\linewidth}
\raggedright
\raisebox{-0.5\height}{\msftbrandmark}
\end{minipage}
\begin{minipage}[c]{0.49\linewidth}
\raggedleft
{\msftdatefont\small\color{msftgray}Month Year}    % <-- Replace with your date
\end{minipage}\par
\vspace{0.35em}
\noindent{\color{msftline}\rule{\linewidth}{0.8pt}\par}

% ============================================================
%  Title + Authors
% ============================================================
\vspace{1.0em}
\begin{center}
{{\msfttitlefont\fontsize{21}{25}\selectfont\color{msftdark}
Your Paper Title Here\par}}                         % <-- Replace with your title
\vspace{1.25em}

{\normalsize\rmfamily\color{msftdark}
Author One$^{1,2,*\dagger}$ \hspace{0.9em}
Author Two$^{1,*}$ \hspace{0.9em}
Author Three$^{2}$ \hspace{0.9em}
Author Four$^{3, \ddagger}$\par
% Usage:
%   $^{number}$   = affiliation
%   $^{*}$        = equal contribution
%   $^{\dagger}$  = internship or special note
%   $^{\ddagger}$ = corresponding author
% Add \\[-0.1em] between lines if you have many authors
}
\vspace{0.22cm}

{\footnotesize\rmfamily\color{msftgray}
$^{1}$ Affiliation One \quad
$^{2}$ Affiliation Two \quad
$^{3}$ Affiliation Three\par
}
\end{center}

% ============================================================
%  Abstract box
% ============================================================
\vspace{0.45em}
\begin{msfttitlebox}
\setlength{\parindent}{0cm}
\setlength{\parskip}{0.14cm}
\raggedright
\nohyphens

% === Abstract ===
% Write your abstract below.
\textbf{Abstract.}
[Background \& Motivation] Provide one or two sentences introducing the broader research area and why the problem matters.
[Problem] Clearly state the specific problem or challenge your work addresses, and briefly mention the limitations of existing approaches.
[Method] Describe your proposed approach at a high level --- what is the key idea or insight that distinguishes your method from prior work?
[Results] Summarize the main experimental findings, including the benchmarks used and the quantitative improvements achieved over baselines.
[Impact] Conclude with the broader implications of your work and potential applications.


\vspace{0.14cm}
{\setlength{\parskip}{0.06cm}\small
% Uncomment and fill in the metadata fields you need:
{\msftmetalabel{Project Page}\href{https://your-project-url}{https://your-project-url}\par}
{\msftmetalabel{Correspondence}
\href{mailto:author3@affiliation.com}{author3@affiliation.com}\par}
{\msftmetalabel{Date}Month Year\par}                % <-- Replace with your date
}
\vspace{0.08cm}
{\footnotesize\rmfamily\itshape\color{msftgray}
$^*$ Equal contribution. \quad
$^{\dagger}$ Work was done during an internship at MSRA. \quad
$^{\ddagger}$ Corresponding author.\par
}
\end{msfttitlebox}
\suppressfloats[t]  % Prevent floats from appearing above the abstract on page 1

% ============================================================
%  Optional teaser figure (uncomment and replace)
% ============================================================
% \begin{figure}[t]
%   \centering
%   \includegraphics[width=\textwidth]{figures/teaser.pdf}
%   \caption{Your teaser figure caption.}
%   \label{fig:teaser}
% \end{figure}

% ============================================================
%  Main content
% ============================================================
% === Introduction ===
\section{Introduction}
\label{sec:introduction}

Your introduction goes here. Describe the problem, motivation, and contributions of your work.

% --- Example: how to cite references ---
For example, the Transformer architecture~\cite{vaswani2017attention} has been widely adopted in natural language processing, computer vision, and multimodal learning. Its self-attention mechanism enables modeling long-range dependencies, making it the backbone of many state-of-the-art systems. Deep residual networks~\cite{he2016deep} enabled training of very deep models by introducing skip connections, which alleviate the vanishing gradient problem and have become a standard component in modern architectures. You can also cite multiple works at once~\cite{devlin2019bert,goodfellow2014generative,kingma2014adam}.

Despite these advances, several challenges remain. [Describe the key challenges or gaps in the current literature that motivate your work. Explain why existing methods are insufficient and what specific limitations you aim to address.]

To address these challenges, we propose \method, a [brief one-sentence description of your approach]. The key insight behind our method is [describe your core idea or observation]. Unlike prior approaches that [describe what existing methods do], our method [describe what makes your approach different and better].

% --- Example: how to reference a figure in Introduction ---
An overview of the proposed framework is illustrated in \cref{fig:method_overview}.

The main contributions of this paper are summarized as follows:
\begin{itemize}[leftmargin=2em,itemsep=2pt,topsep=2pt]
  \item We propose \method, a novel framework for [your task].
  \item We introduce [your key technique or insight], which enables [describe the benefit].
  \item Extensive experiments on [your benchmarks] demonstrate that \method achieves state-of-the-art performance, outperforming existing methods by [X\%] on [metric].
\end{itemize}

% === Related Work ===
\section{Related Work}
\label{sec:related_work}

% --- Example: how to organize related work with bold paragraph headings ---
\noindent\textbf{Topic A.}
Discuss the first line of related work here. Summarize key prior methods, their strengths, and their limitations compared to your approach~\cite{vaswani2017attention,he2016deep}.

\noindent\textbf{Topic B.}
Discuss the second line of related work here. Explain how these methods relate to your problem setting and what gaps remain~\cite{devlin2019bert,goodfellow2014generative}.

\noindent\textbf{Topic C.}
Discuss the third line of related work here. Highlight what distinguishes your approach from existing solutions~\cite{kingma2014adam}.

% === Method ===
\section{Method}
\label{sec:method}

In this section, we describe our proposed method. We first formulate the problem in \cref{sec:formulation}, then present the overall architecture in \cref{sec:architecture}, and finally describe the training procedure in \cref{sec:training}.

% --- Example: how to include a figure ---
\begin{figure}[t]
  \centering
  \fbox{\rule{0pt}{2in}\rule{0.9\linewidth}{0pt}}  % placeholder box; replace with:
  % \includegraphics[width=\linewidth]{figures/your_figure.pdf}
  \caption{Overview of the proposed \method framework. Replace this placeholder with your method pipeline figure.}
  \label{fig:method_overview}
\end{figure}

% --- Example: how to write an equation ---
\subsection{Problem Formulation}
\label{sec:formulation}

Given a dataset $\mathcal{D} = \{(x_i, y_i)\}_{i=1}^{N}$, where $x_i$ denotes the input and $y_i$ denotes the label, our goal is to learn a mapping $f: \mathcal{X} \rightarrow \mathcal{Y}$ that minimizes the expected loss:
\begin{equation}
  \mathcal{L} = \mathbb{E}_{(x,y) \sim \mathcal{D}} \left[ \ell\bigl(f(x; \theta),\, y\bigr) \right]
  \label{eq:objective}
\end{equation}
where $\theta$ denotes the model parameters and $\ell(\cdot, \cdot)$ is the loss function. We refer to \cref{eq:objective} as the training objective throughout this paper.

% --- Example: multi-line equation with a SINGLE number ---
% Use "aligned" inside "equation" for one number; use "align" for per-line numbers.
\begin{equation}
  \begin{aligned}
    \nabla_\theta \mathcal{L}
    &= \frac{1}{N} \sum_{i=1}^{N} \nabla_\theta \ell\bigl(f(x_i; \theta),\, y_i\bigr) \\
    &\approx \frac{1}{|\mathcal{B}|} \sum_{i \in \mathcal{B}} \nabla_\theta \ell\bigl(f(x_i; \theta),\, y_i\bigr),
  \end{aligned}
  \label{eq:gradient}
\end{equation}
where $\mathcal{B}$ denotes a mini-batch sampled from $\mathcal{D}$.

\subsection{Architecture}
\label{sec:architecture}

Describe the architecture of your model here. Explain each component and how they interact with each other.

% --- Example: pseudocode ---
\begin{algorithm}[t]
  \caption{Training Procedure of \method}
  \label{alg:training}
  \KwIn{Dataset $\mathcal{D}$, learning rate $\eta$, number of epochs $T$}
  \KwOut{Trained model parameters $\theta$}
  Initialize $\theta$ randomly\;
  \For{$t = 1$ \KwTo $T$}{
    Sample mini-batch $\mathcal{B} \subset \mathcal{D}$\;
    Compute loss $\mathcal{L}(\theta; \mathcal{B})$ via \cref{eq:objective}\;
    Compute gradient $g \leftarrow \nabla_\theta \mathcal{L}$\;
    Update $\theta \leftarrow \theta - \eta \cdot g$\;
    \If{convergence criterion met}{
      \textbf{break}\;
    }
  }
  \Return{$\theta$}
\end{algorithm}

The training procedure of \method is summarized in \cref{alg:training}.

\subsection{Training Procedure}
\label{sec:training}

Describe how the model is trained, including the optimizer, learning rate schedule, data augmentation, and any other relevant details.

% === Experiments ===
\section{Experiments}
\label{sec:experiments}

We evaluate \method on [your benchmarks] and compare against state-of-the-art baselines. We first describe the experimental setup in \cref{sec:setup}, then present the main results in \cref{sec:main_results}, and finally provide ablation studies in \cref{sec:ablation}.

\subsection{Experimental Setup}
\label{sec:setup}

\noindent\textbf{Datasets.} Describe the datasets used for evaluation, including their size, splits, and any preprocessing steps.

\noindent\textbf{Baselines.} List the baseline methods you compare against and briefly describe their key characteristics.

\noindent\textbf{Implementation Details.} Provide training details such as batch size, learning rate, number of epochs, hardware, and framework used.

\subsection{Main Results}
\label{sec:main_results}

% --- Example: how to include a table ---
\begin{table}[t]
  \centering
  \caption{Main results on [your benchmark]. Best results are in \textbf{bold}. Replace with your own results.}
  \label{tab:main_results}
  \begin{tabular}{lcc}
    \toprule
    Method & Metric A ($\uparrow$) & Metric B ($\downarrow$) \\
    \midrule
    Baseline A & 75.3 & 0.42 \\
    Baseline B & 78.6 & 0.38 \\
    Baseline C & 80.2 & 0.35 \\
    \midrule
    \method (Ours) & \textbf{82.1} & \textbf{0.31} \\
    \bottomrule
  \end{tabular}
\end{table}

As shown in \cref{tab:main_results}, our method outperforms all baselines on both metrics. Specifically, \method achieves a Metric A score of 82.1, surpassing the best baseline by 1.9 points.

% --- Example: how to include subfigures ---
\begin{figure}[t]
  \centering
  \begin{subfigure}[b]{0.48\linewidth}
    \centering
    \fbox{\rule{0pt}{1.2in}\rule{0.8\linewidth}{0pt}}
    \caption{Qualitative result A}
  \end{subfigure}
  \hfill
  \begin{subfigure}[b]{0.48\linewidth}
    \centering
    \fbox{\rule{0pt}{1.2in}\rule{0.8\linewidth}{0pt}}
    \caption{Qualitative result B}
  \end{subfigure}
  \caption{Qualitative comparison. Replace with your visualization results.}
  \label{fig:qualitative}
\end{figure}

\cref{fig:qualitative} shows qualitative comparisons between our method and baselines. Our method produces more [accurate/realistic/consistent] results.

\subsection{Ablation Study}
\label{sec:ablation}

% --- Example: ablation table ---
\begin{table}[h!]
  \centering
  \caption{Ablation study on key components of \method.}
  \label{tab:ablation}
  \begin{tabular}{ccc}
    \toprule
    Component A & Component B & Metric A ($\uparrow$) \\
    \midrule
                &             & 75.3 \\
    \checkmark  &             & 79.4 \\
                & \checkmark  & 78.1 \\
    \checkmark  & \checkmark  & \textbf{82.1} \\
    \bottomrule
  \end{tabular}
\end{table}

We conduct ablation experiments to validate the effectiveness of each component, as shown in \cref{tab:ablation}. Both components contribute to the final performance.

% === Conclusion ===
\section{Conclusion}
\label{sec:conclusion}

In this paper, we presented \method, a novel approach for [your task]. Our method addresses [the key challenge] by [your key insight]. Extensive experiments on [your benchmarks] demonstrate that \method achieves state-of-the-art performance while maintaining [efficiency/simplicity/generality].

\noindent\textbf{Limitations and Future Work.}
Discuss the limitations of your approach and potential directions for future research here.


% ============================================================
%  Bibliography
% ============================================================
\FloatBarrier  % Ensure all floats appear before References
\bibliography{reference}
\bibliographystyle{unsrtnat}

% ============================================================
%  Appendix
% ============================================================
% === Appendix ===
\newpage
\appendix
\setcounter{table}{0}
\setcounter{figure}{0}
\setcounter{section}{0}
\renewcommand{\thetable}{\Alph{section}\arabic{table}}
\renewcommand{\thefigure}{\Alph{section}\arabic{figure}}
\renewcommand{\thesection}{\Alph{section}}

% --- Appendix Table of Contents (toggle via \appendixtoctrue / \appendixtocfalse in main.tex) ---
\ifappendixtoc
  \startcontents[appendix]
  \noindent{\LARGE\bfseries Appendix}
  \vspace{1.5em}
  \titlecontents{lsection}[0em]{\large\bfseries\color{msftblue}\vspace{0.8em}}{\thecontentslabel\hspace{0.5em}}{}{\titlerule*[0.5pc]{.}\contentspage}
  \titlecontents{lsubsection}[2.5em]{\normalsize\color{msftblue}\vspace{0.3em}}{\thecontentslabel\hspace{0.5em}}{}{\titlerule*[0.5pc]{.}\contentspage}
  \printcontents[appendix]{l}{1}{\setcounter{tocdepth}{2}}
  \newpage
\fi

% ============================================================
%  Appendix A: Additional Results
% ============================================================
\section{Additional Results}
\label{sec:additional_results}

This section provides supplementary experimental results that could not be included in the main paper due to space constraints.

\subsection{Extended Comparison}

% --- Example: additional results table in appendix ---
\begin{table}[h!]
  \centering
  \caption{Extended comparison with additional baselines on [your benchmark].}
  \label{tab:extended_results}
  \begin{tabular}{lccc}
    \toprule
    Method & Metric A ($\uparrow$) & Metric B ($\downarrow$) & Metric C ($\uparrow$) \\
    \midrule
    Baseline D & 72.1 & 0.48 & 65.3 \\
    Baseline E & 74.8 & 0.44 & 68.7 \\
    Baseline F & 77.3 & 0.39 & 71.2 \\
    \midrule
    \method (Ours) & \textbf{82.1} & \textbf{0.31} & \textbf{76.9} \\
    \bottomrule
  \end{tabular}
\end{table}

\subsection{Additional Visualizations}

% --- Example: appendix figure ---
\begin{figure}[h!]
  \centering
  \fbox{\rule{0pt}{1.5in}\rule{0.9\linewidth}{0pt}}
  \caption{Additional qualitative results. Replace with your supplementary visualizations.}
  \label{fig:additional_vis}
\end{figure}

\subsection{Implementation Details}

Provide additional implementation details here, such as hyperparameter settings, data preprocessing pipelines, or hardware specifications that are too detailed for the main text.

% ============================================================
%  Appendix B: Prompt Templates
% ============================================================
\section{Prompt Templates}
\label{sec:prompt_templates}

This section documents the prompts used in our pipeline.

\subsection{Evaluation Prompt}

% --- Example: system prompt with natural language ---
\begin{promptbox}[System Prompt for Evaluation]
\noindent\textbf{Role:} You are a professional evaluation expert.

\noindent\textbf{Task:} Given an image and a text description, evaluate whether the image faithfully reflects the content described in the text.

\noindent\textbf{Criteria:}
\begin{itemize}[leftmargin=1.5em,itemsep=1pt,topsep=2pt]
  \item \textbf{Relevance}: Does the image match the text description?
  \item \textbf{Quality}: Is the image free of artifacts and distortions?
  \item \textbf{Consistency}: Are all elements spatially coherent?
\end{itemize}

\noindent\textbf{Output:} Provide a score from 0 to 10 with a brief justification.
\end{promptbox}

\subsection{Structured Output Prompt}

% --- Example: prompt with JSON output format ---
\begin{promptbox}[System Prompt with JSON Output]
\noindent You are given a research paper abstract. Extract the following metadata and return it as a JSON object.

\vspace{0.5em}
\noindent\textbf{Expected output format:}
\begin{lstlisting}[style=json]
{
  "title": "Paper title",
  "authors": ["Author One", "Author Two"],
  "keywords": ["keyword1", "keyword2", "keyword3"],
  "task_type": "classification | generation | detection",
  "datasets": [
    {"name": "ImageNet", "split": "val", "size": 50000}
  ],
  "metrics": {
    "accuracy": 82.1,
    "f1_score": 0.79
  }
}
\end{lstlisting}

\noindent Return \textbf{only} the JSON object, with no additional text.
\end{promptbox}

% ============================================================
%  Appendix C: Theoretical Analysis (optional)
% ============================================================
\section{Theoretical Analysis}
\label{sec:theoretical_analysis}

% This section is optional. Use it when your work involves theoretical contributions.
% The template supports colored theorem boxes: Theorem (blue), Proposition (green),
% Lemma (blue), Corollary (green), Definition (yellow), Assumption (yellow), Remark (gray).
%
% Usage:  \begin{theorem}{Title}{label}  ...  \end{theorem}
% Refer:  \cref{thm:label}, \cref{prop:label}, \cref{lem:label}, etc.

\begin{theorem}{Example Theorem}{example}
Let $f: \mathcal{X} \rightarrow \mathbb{R}$ be an $L$-Lipschitz continuous function. Then for any $x, y \in \mathcal{X}$:
\begin{equation}
  |f(x) - f(y)| \leq L \|x - y\|.
\end{equation}
\end{theorem}

\begin{proof}
Your proof goes here. We omit the details for brevity.
\end{proof}


\end{document}
